\src{189}
par exemple résultant de la commutativité du diagramme suivant d'homomorphismes
canoniques :
\begin{equation}\tag{7.6}
\resizebox{\textwidth}{!}{$
\left\{
\begin{array}{ccccc}
\symup{\operatorname{Ext}_{\Osh_X}^{n-m}(X;\underline{\Omega}_X^{m},\underline{\Omega}_X^{n})}
&
\symup{\longleftarrow}
&
\symup{\operatorname{Ext}_{\Osh_X}^{n-m}(X;\underline{\Omega}_Y^{m},\underline{\Omega}_X^{n})}
&
\symup{\xrightarrow{\sim}}
&
\symup{\operatorname{Hom}_{\Osh_X}(X;\underline{\Omega}_Y^{m},\underline{\Omega}_Y^{m})}
\\[3pt]
\symup{\Big\downarrow}
&&&&
\symup{\Big\downarrow}
\\[3pt]
\symup{\operatorname{Ext}_{\Osh_X}^{n-m}(X;\Osh_X,\underline{\Omega}_X^{n-m})}
&
\symup{\longleftarrow}
&
\symup{\operatorname{Ext}_{\Osh_X}^{n-m}(X;\Osh_Y,\underline{\Omega}_X^{n-m})}
&
\symup{\xrightarrow{\sim}}
&
\symup{\operatorname{Hom}_{\Osh_X}(X;\underline{\Omega}_X^{m},\underline{\Omega}_Y^{m})}
\end{array}
\right.
$}
\end{equation}

On obtient ainsi un équivalent exact du formalisme de la dualité de Poincaré
sur les variétés orientées compactes. En particulier, le théorème 4 permet de
déterminer la classe de cohomologie associée à la diagonale de $\symup{X\times X}$. Par
un raisonnement bien connu, on en déduit par exemple une \uline{formule de Lefschetz} :

\result{THÉORÈME}{5} \uline{Soit} $\symup{f}$ \uline{un endomorphisme d'une variété projective non
singulière} $\symup{X}$, \uline{telle que les points fixes de} $\symup{f}$ \uline{soient de multiplicité} $\symup{1}$.
\uline{Alors le nombre de ces points fixes est congru}, $\symup{\operatorname{mod}}$ \uline{la caractéristique de}
$\symup{\underline{k}}$ \uline{à la somme alternée des traces des endomorphismes des} $\symup{H^{i}(X)}$ \uline{définis
par} $\symup{f}$.

La restriction sur $\symup{f}$ qu'on a dû faire tient aux difficultés mentionnées dans
la remarque du n\textsuperscript{o} 4. On notera, toutefois, que la formule de
Lefschetz reste vraie si $\symup{f}$ est \fquote{homotope} à un endomorphisme dont tous les
points fixes sont de multiplicité $\symup{1}$.

\fsection{\normalfont8}{\normalfont\uline{Généralisation du théorème de dualité}.}

Soit $\symup{X}$ une variété algébrique non singulière telle que tout faisceau
algébrique cohérent $\symup{F}$ sur $\symup{X}$ soit isomorphe à un quotient d'un faisceau
algébrique cohérent localement libre (ce qui est le cas si $\symup{X}$ est localement
fermé dans un espace projectif). Alors tout faisceau algébrique cohérent $\symup{F}$ sur
$\symup{X}$ admet une résolution finie $\symup{L}$ par des faisceaux localement libres, et pour
deux telles résolutions on peut toujours trouver une troisième, et un
homomorphisme de celle-ci dans les deux premières compatibles avec les
augmentations. De même, si $\symup{L}$ est une résolution finie localement libre de
$\symup{F}$, et si on se donne un homomorphisme $\symup{F'\longrightarrow F}$, alors il existe
une résolution finie localement libre $\symup{L'}$ de $\symup{F'}$ et un homomorphisme
$\symup{L'\longrightarrow L}$ compatible avec $\symup{F'\longrightarrow F}$, qu'on peut
supposer surjectif si $\symup{F'\longrightarrow F}$ l'est. Cela permet de définir, étant
donnés deux entiers $\symup{r,s\geq0}$, deux multifoncteurs, cohomologiques, en des
arguments $\symup{A_1,\ldots,A_r}$ ; $\symup{B_1,\ldots,B_s}$

